%%%%%%%%%%%%%%%%%%%%%%%%%%%%%%%%%%%%%%%%%%%%%%%%%%%%%%%%%%%%%%%%%%%%%%%%%%%
% Handout und Uebungsblatt zu Lehreinheit 13
% Im Skript am Kapitelende, in handouts.tex als Druckvorlage.
%%%%%%%%%%%%%%%%%%%%%%%%%%%%%%%%%%%%%%%%%%%%%%%%%%%%%%%%%%%%%%%%%%%%%%%%%%%
\begin{handout}{Wissenschaftsethik und Grenzarbeit}

\textbf{Worum es heute ging.} Ein Trick wurde gezeigt und nicht erklärt.
Die Methode ist kein großes Geheimnis, das Erleben ist es: Eine Enthüllung
hätte die Erfahrung verwandelt, nicht ergänzt. Die Wissenschaft steht auf
der anderen Seite dieser Grenze. Sie muss ihr Verfahren offenlegen -- und
zugleich die Menschen schützen, die an ihr teilnehmen.

\textbf{Kernbegriffe}
\begin{description}
  \item[Grenzarbeit] das Ziehen, Verteidigen und Verschieben von Grenzen
    zwischen Wissen und Nichtwissen, Zeigen und Verbergen
    \parencite{gieryn1983boundary}.
  \item[Informierte Einwilligung] freiwillige, widerrufbare Zustimmung
    nach verständlicher Aufklärung über Zweck, Ablauf, Risiken und Daten.
  \item[Pseudonymisiert / anonym] pseudonymisierte Daten sind über eine
    getrennte Codeliste zuordenbar und bleiben personenbezogen; anonyme
    Daten nicht mehr \parencite{dsgvo2016}.
  \item[FFP] Fabrication, Falsification, Plagiarism: die drei Kerndelikte
    wissenschaftlichen Fehlverhaltens.
  \item[Präregistrierung] Hypothesen und Auswertungsplan vor der
    Datenerhebung festlegen \parencite{nosek2018preregistration}.
\end{description}

\textbf{Offenlegen oder schützen?}

\begin{center}\small
\begin{tabular}{p{0.44\textwidth}p{0.44\textwidth}}
\toprule
\textbf{muss offenliegen} & \textbf{darf/muss geschützt bleiben} \\
\midrule
Methode, Material, Auswertung & Identität der Teilnehmenden \\
Daten, soweit möglich (FAIR) & sensible Rohdaten \\
Finanzierung, Interessenkonflikte & Vertraulichkeit von Gutachten \\
Abweichungen vom Plan & Rechte Dritter \\
\bottomrule
\end{tabular}
\end{center}

\textbf{DSGVO-Grundsätze (Art.~5) in einem Satz.} Nur mit Rechtsgrundlage
und offen, nur für den genannten Zweck, nur so viel wie nötig, korrekt,
nicht länger als nötig, sicher gespeichert -- und nachweisbar so gemacht.

\textbf{Gute wissenschaftliche Praxis.} Der Kodex der DFG beschreibt sie in
19 Leitlinien, u.\,a.\ zu Dokumentation, Archivierung, Autorschaft und
öffentlichem Zugang zu Ergebnissen \parencite{dfg2019kodex}. Bei Verdacht
auf Fehlverhalten helfen die Ombudspersonen Ihrer Hochschule.

\begin{merkbox}
Der Zauberkünstler schützt seine Methode, die Forschung schützt ihre
Teilnehmenden. Offenlegen muss sie alles, was andere zur Prüfung brauchen.
\end{merkbox}

\textbf{Zum Weiterlesen:} \textcite{dfg2019kodex};
\textcite{merton1973normative}; \textcite{dufourmantelle2021geheimnis}.
\end{handout}

\begin{uebungsblatt}{Ethik im eigenen Projekt}

\aufgabe{Ethik-Checkliste für Ihr Semesterprojekt}[in der Stunde, Einzelarbeit, 15 Min.; fertigstellen als Hausaufgabe]
Kreuzen Sie an und notieren Sie für jedes \enquote{nein} oder
\enquote{offen} eine Maßnahme.

\begin{center}\small
\renewcommand{\arraystretch}{1.6}
\begin{tabular}{|p{0.52\textwidth}|c|c|c|p{0.2\textwidth}|}
\hline
\textbf{Frage} & \textbf{ja} & \textbf{nein} & \textbf{offen} & \textbf{Maßnahme} \\
\hline
\multicolumn{5}{|l|}{\textit{Teilnehmende und Einwilligung}} \\ \hline
Ich arbeite mit Menschen (Befragung, Beobachtung, Test). & & & & \\ \hline
Alle erhalten vorher eine verständliche Information. & & & & \\ \hline
Teilnahme ist freiwillig und jederzeit abbrechbar. & & & & \\ \hline
Es nehmen Minderjährige oder Abhängige teil (z.\,B.\ eigene Schüler). & & & & \\ \hline
Ich täusche Teilnehmende; ein Debriefing ist geplant. & & & & \\ \hline
\multicolumn{5}{|l|}{\textit{Daten}} \\ \hline
Ich erhebe nur, was ich für die Frage brauche. & & & & \\ \hline
Namen/Kontaktdaten stehen nicht im Datensatz (Code, getrennte Liste). & & & & \\ \hline
Ich erhebe sensible Daten (Gesundheit, Religion, \ldots). & & & & \\ \hline
Speicherort und Löschtermin sind festgelegt. & & & & \\ \hline
\multicolumn{5}{|l|}{\textit{Integrität und Offenheit}} \\ \hline
Hypothesen und Auswertung sind vor der Erhebung notiert. & & & & \\ \hline
Ausschlüsse von Daten sind begründet und werden berichtet. & & & & \\ \hline
Alle Quellen sind belegt; KI-Nutzung ist nach den Regeln gekennzeichnet. & & & & \\ \hline
Material und Daten (anonymisiert) kann ich anderen zeigen. & & & & \\ \hline
Ich habe keine Interessenkonflikte -- oder lege sie offen. & & & & \\ \hline
\end{tabular}
\end{center}

\aufgabe{Einwilligung: Mustertext anpassen}[in der Stunde, 10 Min.]
Der folgende Text ist ein \emph{Muster} für ein studentisches Projekt ohne
sensible Daten. Er ersetzt keine Prüfung durch die zuständigen Stellen
Ihrer Hochschule. Füllen Sie die Lücken für Ihr Projekt aus und streichen
Sie, was nicht zutrifft.

\begin{center}
\fbox{\parbox{0.92\textwidth}{\small
\textbf{Information und Einwilligung zur Teilnahme}\par\smallskip
Im Rahmen der Lehrveranstaltung \enquote{Wissenschaftliches Arbeiten} an
der \rule{40mm}{0.3pt} untersuche ich, \rule{70mm}{0.3pt}. Die Teilnahme
besteht aus \rule{60mm}{0.3pt} und dauert etwa \rule{10mm}{0.3pt} Minuten.
Es sind keine Risiken bekannt, die über den Alltag hinausgehen.\par\smallskip
Die Teilnahme ist freiwillig. Sie können jederzeit und ohne Angabe von
Gründen abbrechen; daraus entstehen Ihnen keine Nachteile. Ich erhebe keine
Namen. Ihre Angaben werden \rule{30mm}{0.3pt} (anonym / unter einem Code)
gespeichert, nur für dieses Projekt ausgewertet, nur in zusammengefasster
Form berichtet und spätestens am \rule{25mm}{0.3pt} gelöscht. Eine einmal
erteilte Einwilligung können Sie jederzeit widerrufen; bereits anonym
ausgewertete Angaben lassen sich danach nicht mehr einzeln entfernen.
\par\smallskip
Fragen und Widerruf: \rule{50mm}{0.3pt} (Name, E-Mail)\par\medskip
$\square$ Ich habe die Information gelesen und bin mit der Teilnahme
einverstanden.\par\medskip
Datum, Unterschrift: \rule{60mm}{0.3pt}}}
\end{center}

\aufgabe{Fallvignetten: FFP oder nicht?}[in der Stunde, Partnerarbeit, 10 Min.]
Alle Fälle sind \emph{erfunden}. Ordnen Sie zu: \textbf{F} Fabrication,
\textbf{V} Falsification (Verfälschung), \textbf{P} Plagiat,
\textbf{Q} fragwürdige Praxis, \textbf{E} anderes ethisches Problem,
\textbf{OK} in Ordnung. Begründen Sie kurz.

\begin{center}\small
\renewcommand{\arraystretch}{1.5}
\begin{tabular}{|c|p{0.62\textwidth}|p{0.06\textwidth}|p{0.14\textwidth}|}
\hline
 & \textbf{Fall} & \textbf{Typ} & \textbf{Begründung} \\
\hline
1 & Lena fehlen sechs Fragebögen. Sie füllt sie selbst aus, \enquote{so, wie die anderen auch geantwortet haben}. & & \\ \hline
2 & Tom nimmt drei Personen aus der Auswertung, weil ihre Werte \enquote{nicht ins Bild passen}. Im Bericht steht davon nichts. & & \\ \hline
3 & Ein Team rechnet fünf Tests, berichtet nur den signifikanten und formuliert die Hypothese so, als hätte sie vorher festgestanden. & & \\ \hline
4 & Sara übernimmt Gliederung und Argumentationsfolge einer im Netz gefundenen Abschlussarbeit, formuliert alle Sätze neu und nennt die Arbeit nicht. & & \\ \hline
5 & Jonas schließt zwei Ausreißer aus, genau nach der Regel in seiner Präregistrierung, und berichtet das. & & \\ \hline
6 & Ein Student speichert die Liste mit Namen und Codes zusammen mit den Antworten in einer Online-Tabelle, die jeder mit dem Link öffnen kann. & & \\ \hline
7 & Eine Gruppe bewertet eine Lern-App. Ein Gruppenmitglied arbeitet als Werkstudent beim Hersteller; im Bericht steht das nicht. & & \\ \hline
8 & Weil das aktuelle Foto unscharf ist, setzt ein Doktorand ein Bild aus einem älteren Versuch ein und beschriftet es als aktuell. & & \\ \hline
9 & In einer Studie mit einem Zaubertrick erfahren die Teilnehmenden vorher nur, es gehe um \enquote{Wahrnehmung}; danach werden sie aufgeklärt und können ihre Daten zurückziehen. & & \\ \hline
\end{tabular}
\end{center}

\end{uebungsblatt}

\begin{uebungsblattb}{4}{Verantwortung und Rückblick}

\aufgabe{Wann wird Intransparenz zur Täuschung?}[in der Stunde, 65--80]
Der Zauberkünstler täuscht und sagt es; der Sozialpsychologe Diederik
Stapel, der Daten für zahlreiche Studien erfand, täuschte und verschwieg es.
Formulieren Sie eine Regel in einem Satz, wann Nicht-Offenlegen zur
Täuschung wird, und prüfen Sie sie an einem Beispiel aus Medien, Werbung
oder Wissenschaft.
\schreiblinien{4}

\aufgabe{Abschluss-Tagebucheintrag}[Hausaufgabe, letzte Seite des Tagebuchs]
\enquote{Was hat Sie im Laufe dieses Semesters am meisten überrascht --
und was hat das mit Wissenschaft zu tun?} Verwenden Sie mindestens zwei der
13 Kernbegriffe des Semesters: Prediction Error, Operationalisierung,
Suchstrategie, Apophenie/Bias, Falsifikation, experimentelle Kontrolle,
Confirmation Bias, bedingte Wahrscheinlichkeit, Intersubjektivität,
Determinismus, Präzision und Struktur, Credibility Cues,
Forschungsethik.

\begin{lehrhinweis}[Erwartungshorizont]
\textbf{Aufgabe~1:} Wichtig sind konkrete Maßnahmen (\enquote{Codeliste in
getrenntem, passwortgeschütztem Ordner auf dem Hochschullaufwerk}) statt
Absichtserklärungen. Projekte mit Minderjährigen, sensiblen Daten oder
Täuschung vor der Erhebung mit der Lehrperson klären.

\textbf{Aufgabe~2:} Die Studierenden sollen merken, dass sie Zweck, Dauer,
Speicherort und Löschtermin noch nicht festgelegt haben. Hinweis: Echte
Anonymität ist bei kleinen Gruppen schwer; dann ehrlich \enquote{unter
einem Code} schreiben.

\textbf{Aufgabe~3:} 1~=~F. 2~=~V (Verfälschen durch unbegründetes,
verschwiegenes Weglassen). 3~=~Q (p-Hacking und HARKing, LE~10); die
Grenze zur Verfälschung ist diskutierbar, weil der Bericht ein falsches
Bild vom Vorgehen gibt. 4~=~P (Struktur- und Ideenplagiat; neue
Formulierungen ändern daran nichts). 5~=~OK (vorab festgelegt und
berichtet). 6~=~E (Verstoß gegen Vertraulichkeit und Datenminimierung,
Art.~5 DSGVO; kein FFP). 7~=~E (nicht offengelegter Interessenkonflikt).
8~=~V (Bildmanipulation; vgl.\ die Bildduplikate bei
\textcite{bik2016prevalence}). 9~=~OK, sofern die Täuschung nötig ist,
keinen Schaden anrichtet und die Aufklärung wie beschrieben erfolgt.

\textbf{Aufgabe~4:} Eine tragfähige Regel: Nicht-Offenlegen wird zur
Täuschung, wenn jemand den Eindruck erweckt, ein Verfahren sei offen oder
ein Befund geprüft, und andere ihre Entscheidungen auf diesen Eindruck
stützen. Die Zauberkunst erfüllt das nicht, weil das Publikum den Vertrag
kennt.
\end{lehrhinweis}
\end{uebungsblattb}
